\chapter{Levantamento de Requisitos} \label{cap:requisitos}

% TODO: APONTAR PARA O "OBJETIVO", PRESENTE NA INTRODUÇÃO, COMO ALICERCE DOS REQUISITOS
Este capítulo apresenta os requisitos definidos para o sistema de coleta de telemetria veicular desenvolvido neste trabalho, composto por três componentes: um \textit{dongle} OBD-II, um aplicativo móvel e um servidor. O \textit{dongle} é responsável pela aquisição, armazenamento e \textit{upload} dos dados do veículo; o aplicativo permite a interação com o usuário e realiza a intermediação da comunicação entre o \textit{dongle} e o servidor; e o servidor recebe, verifica, armazena e disponibiliza os dados coletados.

O sistema contempla duas formas de utilização. No modo diagnóstico, o aplicativo permite consultar parâmetros e códigos de falha do veículo por meio do \textit{dongle}. No modo coletor, o \textit{dongle} adquire dados automaticamente, encaminhando-os ao servidor por intermédio do aplicativo.

A definição dos requisitos se deu à luz dos objetivos a serem cumpridos pelo projeto (\autoref{sec:intro-objs}) e das preocupações de segurança identificadas na modelagem de ameaças do \autoref{cap:modelagemameacas}. Os requisitos são assim divididos em funcionais, que capturam os comportamentos esperados do sistema, e não funcionais, que abarcam propriedades e restrições que devem ser satisfeitas. Cada requisito receberá um identificador único, no formato RFXY, caso funcional, ou RNFXY, caso não funcional.

\section{Requisitos Funcionais}

Esta seção descreve as funcionalidades esperadas do sistema, agrupadas por fluxo de operação.

\subsection{Diagnóstico e interação com o veículo}

O sistema deve permitir que o usuário consulte informações disponibilizadas pelo veículo por meio da interface OBD-II, utilizando o aplicativo móvel como interface de interação.


\begin{longtable}{p{0.11\textwidth}p{0.79\textwidth}}
\caption{Requisitos funcionais de diagnóstico e interação com o veículo}
\label{tab:rf-diagnostico}\\
\hline
\textbf{ID} & \textbf{Requisito} \\
\hline
\endfirsthead

\hline
\textbf{ID} & \textbf{Requisito} \\
\hline
\endhead

RF01 & O sistema deve permitir a consulta de parâmetros
padronizados do veículo por meio de requisições OBD-II sobre CAN. \\

RF02 & O sistema deve permitir a leitura de códigos de
diagnóstico de falhas (DTCs) disponibilizados pelo veículo. \\

RF03 & O aplicativo deve permitir a visualização dos parâmetros
consultados e dos resultados das operações de diagnóstico. \\

RF04 & O sistema deve permitir a alternância entre os modos
de operação Coletor e Diagnóstico. \\

\hline
\end{longtable}

\subsection{Coleta e Armazenamento de Dados}

Além das consultas ad-hoc realizadas pelo usuário no modo Diagnóstico, o sistema, no modo Coletor, deve realizar a aquisição automática de parâmetros veiculares previamente configurados. A coleta deve ocorrer independentemente da disponibilidade de conexão com o aplicativo ou o servidor.


\begin{longtable}{p{0.11\textwidth}p{0.79\textwidth}}
\caption{Requisitos funcionais de aquisição e armazenamento}
\label{tab:rf-coleta}\\
\hline
\textbf{ID} & \textbf{Requisito} \\
\hline
\endfirsthead

\hline
\textbf{ID} & \textbf{Requisito} \\
\hline
\endhead

RF05 & O sistema deve permitir a configuração do protocolo CAN,
dos PIDs coletados, de seus períodos de aquisição e do prazo
máximo de resposta das requisições. \\

RF06 & No modo Coletor, o \textit{dongle} deve consultar automaticamente
os PIDs configurados, registrando individualmente os valores
obtidos e seus instantes relativos de aquisição. \\

RF07 & O \textit{dongle} deve disponibilizar ao aplicativo os valores
mais recentes dos PIDs coletados, sem exigir novas requisições
ao barramento CAN. \\

RF08 & O sistema deve permitir o recebimento e o registro de
informações de localização fornecidas pelo aplicativo móvel. \\

RF09 & O \textit{dongle} deve armazenar os dados coletados em memória
não volátil, permitindo sua posterior transmissão ao servidor. \\

RF10 & O \textit{dongle} deve organizar os registros em blocos
identificados sequencialmente, preservando as informações
necessárias à recuperação da ordem e da referência temporal
das leituras. \\

\hline
\end{longtable}


\subsection{Comunicação e Sincronização}

A transmissão de telemetria para o servidor depende da comunicação entre os três componentes. O aplicativo atua como intermediário, enquanto o \textit{dongle} e o servidor mantêm os mecanismos necessários para autenticar mensagens e controlar o progresso da sincronização.


\begin{longtable}{p{0.11\textwidth}p{0.79\textwidth}}
\caption{Requisitos funcionais de comunicação e sincronização}
\label{tab:rf-sincronizacao}\\
\hline
\textbf{ID} & \textbf{Requisito} \\
\hline
\endfirsthead

\hline
\textbf{ID} & \textbf{Requisito} \\
\hline
\endhead

RF11 & O sistema deve permitir a comunicação entre \textit{dongle} e
aplicativo por Bluetooth Low Energy (BLE), contemplando
comandos de diagnóstico, configuração e transferência
de dados. \\

RF12 & O aplicativo deve encaminhar ao servidor os blocos
coletados pelo \textit{dongle} e retransmitir os blocos pendentes
em caso de falha de comunicação. \\

RF13 & O sistema deve permitir a autenticação e a sincronização
entre \textit{dongle} e servidor mediante a troca de mensagens
intermediadas pelo aplicativo. \\

RF14 & O servidor deve verificar os blocos recebidos e emitir
confirmações autenticadas dos dados aceitos. \\

RF15 & O \textit{dongle} deve manter os blocos ainda não confirmados
pelo servidor e permitir a retomada de seu envio após
interrupções de conectividade. \\

RF16 & O sistema deve identificar e registrar perdas de blocos
decorrentes de reinicializações, corrupção de dados ou
esgotamento do armazenamento local. \\

RF17 & O sistema deve permitir a obtenção de referências
temporais entre o \textit{dongle} e o servidor, possibilitando
estimar os horários das leituras coletadas. \\

\hline
\end{longtable}


\subsection{Processamento e Consulta de Dados}

Após receber os dados, o servidor deve manter os registros necessários à verificação das informações recebidas e disponibilizar os dados de telemetria para análise.


\begin{longtable}{p{0.11\textwidth}p{0.79\textwidth}}
\caption{Requisitos funcionais de processamento e consulta}
\label{tab:rf-processamento}\\
\hline
\textbf{ID} & \textbf{Requisito} \\
\hline
\endfirsthead

\hline
\textbf{ID} & \textbf{Requisito} \\
\hline
\endhead

RF18 & O servidor deve armazenar os blocos recebidos em seu
formato original e manter as informações necessárias ao
acompanhamento do estado de cada \textit{dongle}. \\

RF19 & O servidor deve decodificar as leituras recebidas
e disponibilizar os dados de telemetria e localização
para consulta e análise. \\

\hline
\end{longtable}


\section{Requisitos Não Funcionais}

Esta seção descreve as propriedades esperadas do sistema nos quesitos de segurança, de confiabilidade e de desempenho. 

\subsection{Segurança}

A principal preocupação de segurança é impedir que o aplicativo intermediário consiga produzir ou modificar registros aceitos como autênticos pelo servidor, alterar sua ordem ou ocultar silenciosamente dados.


\begin{longtable}{
    p{0.10\textwidth}
    p{0.64\textwidth}
    p{0.16\textwidth}
}
\caption{Requisitos não funcionais de segurança}
\label{tab:rnf-seguranca}\\
\hline
\textbf{ID} & \textbf{Requisito} & \textbf{Ameaças} \\
\hline
\endfirsthead

\hline
\textbf{ID} & \textbf{Requisito} & \textbf{Ameaças} \\
\hline
\endhead

RNF01 & O servidor deve verificar criptograficamente a origem
dos blocos de telemetria, aceitando apenas aqueles autenticados
pelo \textit{dongle} correspondente. & A3.3 \\

RNF02 & Alterações no conteúdo autenticado de um bloco devem
ser detectadas, impedindo sua aceitação pelo servidor.
& A3.3 \\

RNF03 & A retransmissão de blocos já aceitos não deve produzir
duplicação dos registros armazenados. & A3.4 \\

RNF04 & O servidor deve preservar a ordem de geração dos
blocos, impedindo que sua transmissão fora de ordem altere
a sequência dos registros aceitos. & A3.5 \\

RNF05 & O sistema deve permitir detectar lacunas na sequência
conhecida de blocos e reconhecer perdas declaradas pelo \textit{dongle},
sem depender da confiabilidade do aplicativo.
& A3.2, A3.6, A3.7 \\

RNF06 & As mensagens de controle entre \textit{dongle} e servidor
devem ter autenticidade e proteção contra reutilização
indevida de desafios e confirmações anteriores.
& A1.2, A3.2 \\

RNF07 & A comunicação entre \textit{dongle} e aplicativo deve utilizar
LE Secure Connections, com mecanismos de pareamento
e restrição de acesso ao dispositivo.
& A1.1, A1.2, A1.4 \\

RNF08 & A comunicação entre aplicativo e servidor deve
utilizar HTTPS/TLS.
& AD4 \\

RNF09 & O \textit{dongle} deve utilizar mecanismos de proteção das
credenciais criptográficas armazenadas e de verificação
da integridade do firmware.
& P1, P3 \\

\hline
\end{longtable}


\subsection{Confiabilidade e Desempenho}

A operação do sistema está sujeita à indisponibilidade de conexão, à capacidade limitada de armazenamento e à variação do tempo de resposta das ECUs. 

\begin{longtable}{p{0.11\textwidth}p{0.79\textwidth}}
\caption{Requisitos não funcionais de confiabilidade e desempenho}
\label{tab:rnf-desempenho}\\
\hline
\textbf{ID} & \textbf{Requisito} \\
\hline
\endfirsthead

\hline
\textbf{ID} & \textbf{Requisito} \\
\hline
\endhead

RNF10 & A aquisição e o armazenamento local da telemetria
devem continuar durante a indisponibilidade do aplicativo
ou do servidor, enquanto o \textit{dongle} estiver operacional
e houver capacidade de armazenamento. \\

RNF11 & A retomada da sincronização após falhas de comunicação
ou reinicializações não deve comprometer a consistência
dos registros anteriormente confirmados pelo servidor. \\

RNF12 & As configurações de aquisição devem respeitar o
limite de utilização de 80\% do canal de requisições,
calculado a partir dos períodos dos PIDs e do prazo
máximo de resposta configurado. \\

RNF13 & Os registros devem preservar a ordem e os intervalos
temporais relativos das leituras dentro de cada inicialização
do \textit{dongle}, mesmo na ausência de uma referência confiável
de tempo absoluto. \\

RNF14 & O aplicativo deve permitir configurar o intervalo
de espera para envio dos blocos ao servidor, possibilitando
ajustar a latência de disponibilização da telemetria. \\

\hline
\end{longtable}


\subsection{Limitações}

Os requisitos não funcionais possuem, como condição de validade, as premissas da modelagem de ameaças (\autoref{sec:premissas}). 

\section{Verificação dos Requisitos}

O atendimento aos requisitos será examinado no \autoref{cap:validacaotestes} por meio de testes funcionais, de desempenho e de segurança. Matrizes de rastreabilidade relacionarão os identificadores aos mecanismos implementados e aos testes executados.